%lezione 13 del 14 aprile 2026
%-----------------------------------------------------------------------------------------------------
\subsection{Esempio: potenziale armonico}
Come sempre, si considera il potenziale armonico come esempio: la sua azione euclidea classica tra
$(q,0)$ e $(q,t)$
è
$$S_{e,cl} = -iS_{cl} = -im\omega q^2\frac{\cos(\omega t)-1}{\sin(\omega t)}
=im\omega q^2\tan\left(\frac{\omega t}{2}\right)\underset{t=-i\tau}
=m\omega q^2\tanh\left(\frac{\omega\tau}{2}\right)$$
Poiché il \textit{path integral} delle deformazioni quadratiche è indipendente da $q$ si calcola
$$\int_{-\infty}^{\infty}dq e^{-\frac{1}{\hbar}S_{e,cl}}
=\int_{-\infty}^{\infty}dq \exp\left\{\frac{-m\omega}{\hbar}\tanh\left(\frac{\omega\tau}{2}\right)
q^2\right\} = \sqrt{\frac{\pi\hbar}{m\omega\tanh\left(\frac{\omega\tau}{2}\right)}}$$
e, separatamente, la correzione quadratica
$$S_{e,2} = \frac{m}{2}\int_{0}^{\tau}d\tau\left(\dv[2]{x}{\tau}+\omega^2x^2\right) = 
\frac{m}{2}\int_{0}^{\tau}d\tau xX_ex$$
con $X_e = -\dv[2]{}{\tau}+\omega^2$.\\
Il \textit{path integral} è, a meno di normalizzazioni, pari a $(\det(X_e))^{-\frac{1}{2}}$:
$$\det(X_e) = \prod_{n=1}^{\infty}\left(\frac{\pi^2n^2}{\tau^2}+\omega^2\right) = 
\frac{1}{\mathcal{N}_e}\prod_{n=1}^{\infty}\left(1+\frac{\omega^2\tau^2}{\pi^2n^2}\right)$$
dove $\mathcal{N}_e$ è la solita normalizzazione della particella libera scritta con il tempo euclideo.

Come in precedenza, il risultato della produttoria si deduce dal fatto che essa, come funzione di
$\omega \tau$, non ha poli e ha zeri in $\omega\tau = \pm i\pi,\pm 2i\pi...$: l'unica funzione
con queste caratteristiche è
$$\prod_{n=1}^{\infty}\left(1+\frac{\omega^2\tau^2}{\pi^2n^2}\right) =
\frac{\sinh(\omega\tau)}{\omega\tau}$$
Si trova dunque che
$$Z = \sqrt{\frac{\pi\hbar}{m\omega\tanh\left(\frac{\omega\tau}{2}\right)}}
\sqrt{\frac{m}{2\pi\hbar\tau}}\sqrt{\frac{\omega\tau}{\sinh(\omega\tau)}} =
\frac{1}{2\sinh\left(\frac{\omega\hbar\beta}{2}\right)}$$
che, come per il propagatore, è il risultato esatto del problema (non essendoci termini dipendenti dalla
terza potenza delle deformazioni). Un altro modo per giungere a questo risultato è calcolare direttamente
la traccia, sfruttando gli autostati di $H$:
$$\Tr(e^{-\beta H}) = \sum_{n=0}^{\infty}\bra{n}e^{-\beta H}\ket{n} =
\sum_{n=0}^{\infty}e^{-\beta\hbar\omega(n+\frac{1}{2})} =
e^{-\frac{\omega\hbar\beta}{2}}\sum_{n=0}^{\infty} e^{-\omega\hbar\beta n} =
\frac{e^{-\frac{\omega\hbar\beta}{2}}}{1-e^{-\omega\hbar\beta}} =
\frac{1}{2\sinh\left(\frac{\omega\hbar\beta}{2}\right)}$$
\textbf{Oss.} In questo caso è stato possibile utilizzare un modo diretto per risolvere il
problema, ma nella maggioranza dei casi non sono noti gli autovalori in energia di $H$ e questo metodo
non è applicabile.
\newpage
\chapter{Approssimazione WKB}\\

L'approssimazione WKB (Wentzel–Kramers–Brillouin) è un metodo di risoluzione per equazioni
differenziali del second'ordine, che qui sarà applicata all'equazione di Schrödinger.
\section{Regime WKB}
Nel capitolo precedente si è considerato il limite semiclassico $\hbar\to0$, qui si cerca un altro
regime: si consideri, senza perdere di generalità, una funzione d'onda scritta in termini di una
funzione complessa $S(\mathbf{x},t)$
$$\psi(\mathbf{x},t) := e^{\frac{i}{\hbar}S(\mathbf{x},t)}$$
Imponendo l'equazione di Schrödinger dipendente dal tempo si trova che
\begin{gather*}
i\hbar\pdv[]{}{t}\psi(\mathbf{x},t) = \frac{-\hbar^2}{2m}\laplacian \psi(\mathbf{x},t) +
V(\mathbf{x})\psi(\mathbf{x},t)\\ \updownarrow \\
-\pdv[]{S}{t} = \frac{1}{2m}\left(\gradient S\right)^2-\frac{i\hbar}{2m}\laplacian S+V(\mathbf{x})
\end{gather*}
che, nel limite $\hbar\to0$, si riduce all'equazione di Hamilton-Jacobi per l'azione classica $S$:
più precisamente, questo limite è equivalente a considerare il modulo del termine proporzionale a
$\hbar$ trascurabile rispetto all'altro termine dipendente da $S$ nel lato destro dell'equazione
$$\left|\hbar\laplacian S\right|<<\left|\gradient S\right|^2\leftrightarrow\frac{\left|\hbar\laplacian 
S\right|}{\left|\gradient S\right|^2}<<1$$
utilizzando la relazione classica $\mathbf{p} = \gradient S$ si trova che ciò vale per
$$\frac{\left|\hbar\divergence \mathbf{p}\right|}{\left|\mathbf{p}\right|^2}<<1$$
che vuol dire
\begin{itemize}
    \item $\left|\div \mathbf{p}\right|$ piccolo, cioè \textbf{impulso lentamente variabile};
    \item $\left|\mathbf{p}\right|^2$ grande, cioè \textbf{impulso grande}.
\end{itemize}
Se si considera ora il caso di $H$ unidimensionale e indipendente dal tempo e si pone\\
$S(x,t) = -Et+W(x)$
$$\psi(x,t) = e^{-\frac{i}{\hbar}S(x,t)} :=e^{-\frac{i}{\hbar}Et}
Ae^{\frac{i}{\hbar}W(x)}:=e^{-\frac{i}{\hbar}Et}\psi(x)$$
si trova, sostituendo nell'equazione di Schrödinger indipendente dal tempo,
$$H\psi = E\psi\leftrightarrow\frac{1}{2m}(W')^2 -\frac{i\hbar}{2m}W''+V(x)-E=0$$
che è detta equazione di Hamilton-Jacobi ridotta.\\
Il suo limite classico si trova, similmente a prima, quando
$$\left|\frac{\hbar W''}{(W')^2}\right|<<1\leftrightarrow\left|\frac{\hbar p'}{p^2}\right|<<1$$
Se si pone $\hbar=0$, dall'equazione si trova
$$(W'(x))^2 = 2m\left[E-V(x)\right]:=p^2(x)\rightarrow p' = \frac{-m}{\sqrt{2m(E-V(x))}}\pdv[]{V}{x}=
\frac{-m}{\pm p}\pdv[]{V}{x}$$
e il limite classico diventa
$$\left|\frac{\hbar p'}{p^2}\right| = \left|\frac{\hbar m\pdv[]{V}{x}}{p^3}\right|<<1$$
che significa \textbf{potenziale lentamente variabile} e \textbf{impulso grande}.
\subsection{Soluzione perturbativa al primo ordine}
Per trovare soluzioni semiclassiche, si risolve perturbativamente, espandendo $W(x)$ in potenze di
$\hbar$:
$$W(x) = W_0(x) +\hbar W_1(x)+\hbar^2W_2(x)+...$$
Imponendo questa forma nell'equazione di Hamilton-Jacobi ridotta ed eguagliando i termini dello stesso
ordine in $\hbar$, all'ordine 0 si ritrova l'equazione classica
$$\frac{1}{2m}(W_0')^2 +V(x)-E=0$$
con soluzione $W_0' :=p(x)\in\mathbb{R}$, mentre al primo ordine si trova
$$\frac{1}{2m}\left(2W_1'W_0'\right)-\frac{i}{2m}W_0'' = 0 \leftrightarrow
W_1'(x) = \frac{iW_0''}{2W_0'}$$
\textbf{Oss.} In tutti gli ordini, si trova un'equazione differenziale lineare rispetto alla funzione
da trovare, con coefficienti dipendenti dagli ordini precedenti (noti), che è sempre risolubile.\\

In questo corso si tratterà solamente dell'approssimazione al primo ordine: la soluzione dell'ordine
0, nel caso in cui $E>V(x)$, è
$$W_0'=\pm\sqrt{2m(E-V(x))}:=p(x)\rightarrow W_0(x)=\pm\int_{x_0}^{x}dx'p(x') =
\pm\int_{x_0}^{x}dx'\sqrt{2m(E-V(x))}$$
dove $x_0$ è un qualsiasi punto in cui vale il regime di approssimazione considerato.\\
Da questa soluzione si trova che
$$W_1'(x) = \frac{i}{2}\frac{p'(x)}{p(x)} = \frac{i}{2}\dv[]{}{x}\ln\left(\frac{p(x)}{p_0}\right) = 
i\dv[]{}{x}\ln\sqrt{\frac{p(x)}{p_0}}\rightarrow W_1(x)=i\ln\sqrt{\frac{p(x)}{p_0}}$$
Sostituendo nell'espressione di $\psi(x)=e^{\frac{i}{\hbar}W(x)}$ si trova
$$\psi(x) = \exp\left\{\pm i\int_{x_0}^{x}dx'\frac{p(x')}{\hbar}\right\}\cdot
\exp\left\{-\ln\sqrt{\frac{p(x)}{p_0}}\right\}$$
infine, ridefinendo $k(x):=\frac{p(x)}{\hbar}$
$$\psi(x) = \frac{C}{\sqrt{k(x)}}\exp\left\{\pm i\int_{x_0}^{x}dx'k(x')\right\}$$
\textbf{Oss.} La densità di proabilità associata a $\psi(x,t)$ è proporzionale a quella classica:
$$\left|\psi(x,t)\right|^2dx=\frac{dx}{k(x)} = \hbar\frac{dx}{p(x)}=m\hbar\frac{dx}{v(x)}=
m\hbar dt\rightarrow|\psi(x,t)|^2\propto \dv[]{t}{x}$$
La soluzione generale dell'equazione di Schrödinger stazionaria è una combinazione lineare delle due
soluzioni linearmente indipendenti
\begin{equation}
\boxed{\psi(x)=\frac{1}{\sqrt{k(x)}}\left(a\exp\left\{i\int_{x_0}^{x}dx'k(x')\right\}+b
\exp\left\{-i\int_{x_0}^{x}dx'k(x')\right\}\right)}
\label{eq:wkb_free_exp}
\end{equation}
che può essere riscritta come
\begin{equation}
    \boxed{\psi(x)=\frac{A}{\sqrt{k(x)}} \sin\left(\int_{x_0}^{x}dx'k(x')+\delta\right)}
\label{eq:wkb_free_sin}
\end{equation}

Nel caso in cui invece $E<V(x)$, si definisce la funzione $\beta(x) = \sqrt{2m(V(x)-E)}$ tale che
$\beta^2(x)=-k^2(x)$ e si giunge, similmente, a
\begin{equation}
    \boxed{\psi(x) = \frac{1}{\sqrt{\beta(x)}}\left(c\exp\left\{\int_{x_0}^{x}dx'\beta(x')\right\} +
    d\exp\left\{-\int_{x_0}^{x}dx'\beta(x')\right\}\right)}
    \label{eq:wkb_constr}
\end{equation}

